\documentclass[11pt]{article}
\usepackage[a4paper,margin=1.1in]{geometry}
\usepackage[T1]{fontenc}
\usepackage{lmodern}
\usepackage{amsmath,amssymb,amsthm}
\usepackage{booktabs}
\usepackage{longtable}
\usepackage{array}
\usepackage{microtype}
\usepackage[hidelinks]{hyperref}
\usepackage{natbib}
\usepackage{setspace}
\usepackage{enumitem}
\usepackage{xcolor}

\newtheorem{theorem}{Theorem}[section]
\newtheorem{proposition}[theorem]{Proposition}
\newtheorem{lemma}[theorem]{Lemma}
\newtheorem{corollary}[theorem]{Corollary}
\theoremstyle{definition}
\newtheorem{definition}[theorem]{Definition}
\newtheorem{assumption}[theorem]{Assumption}
\theoremstyle{remark}
\newtheorem{remark}[theorem]{Remark}

% A formal name, as it appears in the Lean development.
\newcommand{\leanunderscore}{\char`\_\hskip 0pt plus 0.4pt\relax}
\newcommand{\lean}[1]{\begingroup\ttfamily\let\_\leanunderscore #1\endgroup}
% The Lean name attached to a stated result.
\newcommand{\formal}[1]{\par\noindent\begingroup\footnotesize\raggedright\emph{Formal statement:} \lean{#1}.\par\endgroup}
\newcommand{\R}{\mathbb{R}}
\newcommand{\E}{\mathbb{E}}
\newcommand{\abar}{\bar a}

\onehalfspacing
\tolerance=600 \emergencystretch=1.5em

\title{Krusell--Smith Economies in Lean:\\ Existence, Uniqueness, Stationarity and the Aggregate State}
\author{LeanEconomics Org [mostly Claude]\thanks{The Lean development described here is at
\url{https://github.com/LeanEconomics/LeanEconomics}, directory \texttt{LeanEconomics/KrusellSmith}.
It was written mostly by Claude (Anthropic), an AI coding model, working under the direction of
Robert Kirkby (Victoria University of Wellington); every result reported was checked by the Lean
kernel, and every headline theorem depends only on Lean's standard axioms. Companion to the
write-ups on Aiyagari (1994) and on overlapping generations in the same repository.}}
\date{November 2026\\[0.5em]\small Draft, in progress}

\begin{document}
\maketitle

\begin{abstract}
\noindent
This is the record of a programme to formalise heterogeneous-agent economies with aggregate risk in
the style of \citet{krusellsmith1998} in Lean~4, and of what it proves. Sequential equilibrium
exists from every initial state, for logarithmic and CRRA utility and for Cobb--Douglas prices.
Recursive equilibrium exists when the aggregate state is enlarged by the value function or by the
infinite history of aggregate shocks, and a stationary recursive equilibrium with an ergodic
distribution exists in the sense of \citet{dgmm1994}. When the capital share is small, sequential
equilibrium is unique, so a recursive equilibrium exists on the minimal state $(z,\mu)$; the law of
motion contracts, which turns the matched-expectations idea of \citet{lee2026} into explicit error
bounds; and the equilibrium is stationary directly from the sequential equilibrium, with a unique
and ergodic invariant distribution. We then ask when aggregate shock and aggregate capital $(z,K)$
suffice as the state. We characterise the answer exactly: a recursive equilibrium on $(z,K)$
exists if and only if two histories with the same $(z,K)$ have the same next capital (under Krusell
and Smith's employment design). We bound the
error of forecasting with $(z,K)$ by the nonlinearity of saving plus the feedback of prices, and give
an explicit economy with heterogeneous households and Cobb--Douglas prices in which $(z,K)$ is
exactly sufficient: sequential equilibrium is unique, $(z,K)$ is a sufficient state, and the recursive
equilibrium on $(z,K)$ has a unique, ergodic stationary distribution, with persistent aggregate
shocks. A second economy, with a linear technology and i.i.d.\ shocks, has households saving
strictly inside their constraints and a cross-section that keeps its whole history, and $(z,K)$ is
still sufficient. Whether $(z,K)$ is sufficient at the interior, persistent calibrations of Krusell
and Smith, with prices responding to capital, remains open.
\end{abstract}

\tableofcontents

\section{Introduction}\label{sec:intro}

\citet{krusellsmith1998} solve a heterogeneous-agent economy with aggregate risk by supposing that
households forecast prices from the aggregate shock and the mean of the wealth distribution, and
find that this forecast is almost exact. Behind the method lie questions the method does not
answer. Does an equilibrium exist? Is it recursive in some state, and in which? Is the economy
stationary? And when is the aggregate state $(z,K)$ --- shock and capital --- actually enough?

This note records what a Lean formalisation of the model proves about these questions. The
results are:

\begin{enumerate}[leftmargin=*,itemsep=2pt]
\item \textbf{Existence} (Section~\ref{sec:existence}). From every aggregate shock and initial
  distribution there is a sequential equilibrium, by Schauder--Tychonoff on paths of capital over
  the tree of shock histories; with logarithmic and CRRA utility; with Cobb--Douglas prices, the
  capital clamp provably slack under an explicit condition.
\item \textbf{Recursive equilibrium} (Section~\ref{sec:recursive}). With the value function, or
  the infinite history of aggregate shocks, added to the state, a recursive equilibrium exists in
  every economy; a stationary recursive equilibrium with an ergodic distribution exists in the
  sense of \citet{dgmm1994}, conditionally spotless in the sense of \citet{cao2020}. Uniqueness of
  sequential equilibrium implies a recursive equilibrium on the minimal state $(z,\mu)$.
\item \textbf{Uniqueness for small capital shares} (Section~\ref{sec:uniqueness}). For CRRA utility
  and Cobb--Douglas prices there is a capital share below which sequential equilibrium is unique,
  checked from primitives by a contraction in marginal utility. The law of motion then contracts in
  the distribution, which bounds the error of matching on capital \citep{lee2026}.
\item \textbf{Stationarity from the sequential equilibrium} (Section~\ref{sec:stationarity}). The
  equilibrium is a function of the infinite history of shocks, so it inherits stationarity from
  the shocks: the aggregate state of the unique small-share equilibrium has a unique, ergodic
  invariant distribution, with no recursive state space and no randomisation.
\item \textbf{The aggregate state $(z,K)$} (Section~\ref{sec:aggregate}). A recursive equilibrium on
  $(z,K)$ representing the equilibrium exists if and only if histories with the same $(z,K)$ have
  the same future capital and labour; then it is continuous and its stationary law is unique.
  Without idiosyncratic heterogeneity this always holds. The error of
  forecasting next capital from $(z,K)$ is at most the nonlinearity of saving plus the price
  feedback, times the distance between cross-sections.
\item \textbf{An explicit heterogeneous economy} (Section~\ref{sec:corner}). Log utility,
  Cobb--Douglas prices, employed households saving to a cap and unemployed households spending
  everything: sequential equilibrium is unique from every initial state, $(z,K)$ is sufficient, and
  the recursive equilibrium on $(z,K)$ has a unique, ergodic stationary distribution. Aggregate
  shocks have Krusell and Smith's persistence.
\item \textbf{Interior saving} (Section~\ref{sec:interior}). With prices depending on the shock
  only (a linear technology), i.i.d.\ shocks and cash ordered by the shock, $(z,K)$ is sufficient in
  every sequence-state equilibrium; in an explicit economy half of the households save strictly
  inside their constraints after every good shock.
\end{enumerate}

\noindent
Section~\ref{sec:open} lists what is open. Appendix~\ref{app:index} indexes the formal statements.

The prior literature is close in places. \citet{miao2006} proves existence of sequential equilibrium
by Schauder--Tychonoff and a recursive equilibrium with payoffs in the state; \citet{cao2020} proves
existence for the Krusell--Smith specification, a generalised recursive equilibrium, the value
function as extended state and, under uniqueness, the minimal-state recursive equilibrium, and
leaves the minimal-state question without uniqueness open; \citet{kublerschmedders2002} give
finite-agent examples where the wealth distribution is not a sufficient state. Both
\citet{miao2006} and \citet{kublerschmedders2002} treat uniqueness as a condition that cannot in
general be verified from primitives; Section~\ref{sec:uniqueness} verifies it from primitives, though
only for small capital shares. \citet{prohl2025} claims existence and uniqueness of a minimal-state recursive equilibrium by
monotone-operator methods; we found gaps in the proofs as written (the interiority argument in the
existence step, and convexity of the admissible set) and do not rely on it.

\section{The model as formalised}\label{sec:model}

Aggregate shocks $z$ take finitely many values and follow a Markov chain $\pi(z,z')$. Each household
has an idiosyncratic state $e$ in a finite set, moving with probabilities $p(z,z',e,e')$ that may
depend on today's and tomorrow's aggregate shock, and labour efficiency $\epsilon(e)$. Households
hold positions $a-b$ with $a\in[0,\abar]$, where $b\ge0$ is a borrowing limit, discount at
$\beta<1$, and receive the gross return $R(z,K,L)$ on their position and the wage $w(z,K,L)$ per
efficiency unit, where $K$ is mean (shifted) assets and $L$ aggregate labour. Cash on hand is
$R(a-b)+w\epsilon(e)$; saving $a'$ is feasible when
$0\le a'\le\min(\abar,\ R(a-b)+w\epsilon(e)-\underline c)$, with $\underline c\ge 0$ a consumption
floor that labour income net of debt service at the limit always affords. Krusell and Smith's
economy has $b=0$; results that need it say so (the monotone route of
Section~\ref{sec:recursive} and the capital lower bound below, since a higher return hurts
debtors). Period utility is continuous, increasing and strictly concave on
$[\underline c,\infty)$.

Prices are abstract bounded continuous functions; the Cobb--Douglas instance is
$R=1-\delta+\alpha A_zK^{\alpha-1}L^{1-\alpha}$, $w=(1-\alpha)A_zK^\alpha L^{-\alpha}$ with capital
and labour clamped to boxes, so that prices are bounded. Utility unbounded below at zero (logarithmic,
CRRA with $\gamma\ge1$) enters through a positive income floor and a consumption floor proved never
to bind (\lean{KSInada}). A household's problem along a path of prices is a stochastic dynamic
program on the tree of shock histories; its value function is the fixed point of a Blackwell
contraction, and strict concavity gives a unique optimal policy, jointly continuous in assets and the
price path.

\begin{definition}
A \emph{sequential equilibrium} from shock $z_0$ and distribution $\mu_0$ is a path $K$ of capital on
the tree of shock histories such that, at every node, capital equals mean assets of the distribution
generated by optimal saving along the price path of $K$ (with labour the mean efficiency of the
idiosyncratic marginal).
\end{definition}

\noindent
Two features of the formalisation should be kept in mind. The asset cap $\abar$ is a constraint;
for small capital shares it is proved irrelevant (Theorem~\ref{thm:cap}). And the clamping boxes of the Cobb--Douglas instance are part of the price
specification; where it matters we prove that the equilibrium stays inside them, so that prices are
the true marginal products.

\section{Existence of sequential equilibrium}\label{sec:existence}

\begin{theorem}\label{thm:existence}
In every Krusell--Smith economy, from every aggregate shock $z_0$ and every initial distribution
$\mu_0$, a sequential equilibrium exists, with aggregate labour equal to the mean efficiency at
every node.
\end{theorem}
\formal{exists\_sequentialEquilibrium, exists\_sequentialEquilibrium\_labour}

\noindent
The unknown is a path of capital on the tree of shock histories, with values in $[0,\abar]$. These
paths form a compact convex subset of a locally convex space in the product topology. The household
value function is continuous in the path: a depth-weighted sup norm makes the Bellman operator a
contraction with a modulus uniform in the path, and the weights make the tree's distant nodes
negligible. Strict concavity gives a unique policy, continuous in the path by Berge's theorem; the
distributions on the tree are then weakly continuous in the path by induction on depth, and so is
mean capital. Schauder--Tychonoff gives a fixed point.

\begin{theorem}\label{thm:existenceInada}
The same holds with utility defined only on $(0,\infty)$ and marginal utility tending to infinity at
zero --- logarithmic and CRRA utility --- provided labour income is bounded below by a positive
floor; optimal consumption never reaches the consumption floor used in the construction.
\end{theorem}
\formal{KSInada.exists\_sequentialEquilibrium}

\begin{theorem}\label{thm:existenceCD}
With logarithmic utility and Cobb--Douglas prices, a sequential equilibrium exists, and prices are
the true marginal products wherever capital lies in the clamping box. With CRRA utility, an explicit
lower end $\underline K$ of the capital box (depending on $\beta$, $\gamma$, the persistence of the
aggregate chain and the technology) keeps every equilibrium's capital above $\underline K$: the
clamp never binds.
\end{theorem}
\formal{logEconomy\_exists\_equilibrium, crra\_exists\_unclamped\_equilibrium}

\noindent
The lower bound adapts \citet[Lemma 9]{cao2020}: a per-household Euler inequality, aggregated with
Jensen's inequality for the convex function $c\mapsto c^{-\gamma}$, bounds consumption today by
expected consumption at the child nodes, and an induction on the tree pushes capital above
$\underline K$.

\section{Recursive equilibrium}\label{sec:recursive}

\subsection{Uniqueness gives the minimal state}

\begin{theorem}\label{thm:minimal}
If sequential equilibrium is unique from every $(z,\mu)$, then there is a recursive equilibrium on
the minimal state $(z,\mu)$: the distribution next period is a function of today's shock, today's
distribution and tomorrow's shock, weakly continuous in the distribution, and saving is a continuous
function of today's state.
\end{theorem}
\formal{KSModel.isRecursive, continuous\_lawOfMotion}

\noindent
The continuation of a sequential equilibrium below a node is a sequential equilibrium from that
node's state, and the graph of the equilibrium correspondence is closed; with uniqueness the
equilibrium path is a continuous function of $(z,\mu)$ by the closed-graph theorem on a compact
space, and the law of motion is read off it. This is \citet[Proposition 2]{cao2020}.

\subsection{Enlarged states, in every economy}

\begin{theorem}\label{thm:valuestate}
In every economy there is a recursive equilibrium whose state is $(z,\mu,v)$, with $v$ the profile
of household values, and every such equilibrium generates a sequential equilibrium with the same
payoffs.
\end{theorem}
\formal{recursiveEquilibrium\_valueState, generates\_sequentialEquilibrium, payoffEquivalent}

\begin{theorem}\label{thm:history}
In every economy there is a recursive equilibrium whose aggregate state is the infinite history of
aggregate shocks: an assignment $\mu^*$ of a distribution to each history, consistent with optimal
saving at every history, with the history evolving exogenously. This holds in particular with
logarithmic utility and Cobb--Douglas prices.
\end{theorem}
\formal{exists\_seqStateEquilibrium, exists\_historyRecursiveEquilibrium,
logEconomy\_exists\_seqStateEquilibrium}

\noindent
The fixed point is taken in the space of assignments of distributions to infinite histories; no
initial distribution is needed. We call it the \emph{sequence-state equilibrium}. Under uniqueness
of sequential equilibrium it is Markov in $(z,\mu)$.

\begin{theorem}\label{thm:dgmm}
In every economy there is a stationary recursive equilibrium with an ergodic invariant distribution,
on the expectations-augmented state of \citet{dgmm1994}; it can be chosen conditionally spotless in
the sense of \citet[Definition 6]{cao2020}.
\end{theorem}
\formal{exists\_stationaryRecursiveEquilibrium, exists\_conditionallySpotlessEquilibrium}

\noindent
The proof formalises DGMM's Theorem 1.1 and Corollary 1.1 in the abstract (Kakutani--Fan--Glicksberg
for a correspondence of measures, disintegration, Krein--Milman for the ergodic choice) and applies it
to the compact set of equilibrium states.

\subsection{The minimal state without uniqueness}

\begin{theorem}\label{thm:monotone}
If saving rises with the path of aggregate capital, the greatest sequential equilibrium is a
consistent selection, and a recursive equilibrium on $(z,\mu)$ exists without uniqueness. Saving
rises with capital in complementarity economies (a constant wage, a return nondecreasing in capital,
relative risk aversion at most one), and an explicit such economy is given.
\end{theorem}
\formal{SavingMonotone.isRecursive, savingMonotone\_of, example\_exists\_consistentSelection}

\noindent
The condition fails for Cobb--Douglas prices, where higher capital lowers the cash on hand of
wealthy households.

\section{Uniqueness for small capital shares, and matching}\label{sec:uniqueness}

\subsection{A contraction in marginal utility}

Uniqueness is proved by showing that the equilibrium map --- capital path to the mean-saving path it
induces --- is a contraction in a depth-weighted sup norm. The household side is measured in
marginal utility, not in policies: the gap in $u'(c)$ between two price paths propagates with factor
$\beta\bar R$ per period, and the policy gap follows by strong concavity on the range of
consumption. With $\ell$ the Lipschitz constant of prices in capital, the contraction factor is
$\kappa=\lambda(\ell C+A/m)/(1-\lambda\theta)$, explicit in the primitives, where $\theta$ is a bound
on the slope of saving in assets.

\begin{theorem}\label{thm:smallshare}
For CRRA utility ($\gamma>0$) and Cobb--Douglas prices with the clamp, there is $\alpha_0>0$ such
that for every capital share $\alpha<\alpha_0$ sequential equilibrium is unique from every state, so
the recursive equilibrium on $(z,\mu)$ of Theorem~\ref{thm:minimal} exists. This holds for
$0\le\delta<1$ whenever $(\beta(1-\delta))^{1/\gamma}<1-\delta$, using the slope bound
$\theta=\bar R(\beta\bar R)^{1/\gamma}/\underline R$ obtained from the minimal marginal propensity to
consume; in particular for $\delta=0$ unconditionally.
\end{theorem}
\formal{crra\_existsUnique\_sequentialEquilibrium, smallShare\_contraction, minMPC\_smallShare}

\noindent
The proof is analytic: $\kappa$ and the rate of the law of motion are explicit functions of $\alpha$,
continuous at $\alpha=0$ where they reduce to $0$ and to the slope bound. The theorem is
qualitative. At the Krusell--Smith calibration the crude constants give $\kappa$ of order $10^{15}$
at $\alpha=0.36$, and $\kappa<1$ only for capital shares of order $10^{-13}$ (a separate numerical
evaluation, not part of the formal development).

\begin{theorem}\label{thm:sequnique}
Under the same conditions, if employment transitions are bounded below, there is exactly one
sequence-state equilibrium.
\end{theorem}
\formal{minMPC\_smallShare\_existsUnique\_seqStateEquilibrium,
minor\_minMPC\_smallShare\_existsUnique\_seqStateEquilibrium}

\subsection{Fading memory and matched expectations}

\citet{lee2026} argues that matching on aggregate capital along a simulated path is accurate because
the economy forgets its past. In the small-share class this is a theorem. Measure distance between
cross-sections by a two-part distance: Lipschitz test functions in assets, and the employment
marginals.

\begin{theorem}\label{thm:matching}
For small capital shares, with both shock chains bounded below, the equilibrium law of motion
contracts at a rate $r^*<1$ in the two-part distance, even when employment marginals depend on the
history; in every sequence-state equilibrium, two histories agreeing over their last $n$ shocks have
savings within $C_0(r^*)^n$, with $C_0$ explicit.
\end{theorem}
\formal{dist2\_run, minMPC\_smallShare\_matching, smallShare\_matching}

\section{Stationarity directly from the sequential equilibrium}\label{sec:stationarity}

The usual route to a stationary equilibrium passes through a recursive state space, often enlarged
and randomised \citep{dgmm1994}. A sequence-state equilibrium offers a direct one: it is a function
of the infinite history of shocks, and the history evolves exogenously, so if the history is drawn
from a law left invariant by the shock process, the aggregate state $(z,\mu^*)$ has the same law
every period. The equilibrium is a factor of the stationary shock process.

\begin{theorem}\label{thm:factor}
The shock-history process has an invariant ergodic law, unique when the aggregate chain is bounded
below. For every sequence-state equilibrium measurable in the history --- with no uniqueness ---
tomorrow's aggregate state has today's law. Under uniqueness of sequential equilibrium that law is
an invariant distribution of the equilibrium Markov process on $(z,\mu)$, ergodic when the shock law
is.
\end{theorem}
\formal{exists\_invariant\_ergodic, stationary\_unique\_of\_minor, stateLaw\_tomorrow,
isInvariant\_stateLaw}

\begin{theorem}\label{thm:stationary}
For CRRA utility, Cobb--Douglas prices with $0\le\delta<1$ and $(\beta(1-\delta))^{1/\gamma}<1-\delta$,
both shock chains bounded below and a small enough capital share: the sequence-state equilibrium is
unique and continuous in the history; the law of its aggregate state under the stationary shock law
is an invariant and ergodic distribution of the equilibrium Markov process; and every invariant
distribution of that process equals it.
\end{theorem}
\formal{continuous\_seqState, minMPC\_smallShare\_unique\_stationary}

\noindent
Continuity comes from fading memory (Theorem~\ref{thm:matching}) and Stone--Weierstrass density of
asset-Lipschitz test functions. Uniqueness of the invariant distribution is a coupling: the $T$-step
expectation of a test function from any state is within $C(r^*)^T$ of a function of today's shock
alone, invariant distributions have the same shock marginal, and the test functions separate
measures on the Polish aggregate state space. The theorem proves uniqueness of the invariant
distribution, not convergence in law from every starting point.

\section{When is $(z,K)$ a sufficient state?}\label{sec:aggregate}

Households need only prices, and prices depend on the path of capital and labour. So the right
question is not whether $(z,K)$ determines the distribution, but whether it determines that path.

\begin{definition}
A sequence-state equilibrium $\mu^*$ satisfies \emph{$(z,K)$-sufficiency} if any two histories with
the same current shock and the same capital have the same capital and labour at every node of their
future trees.
\end{definition}

\noindent
A \emph{recursive equilibrium on $(z,K)$} consists of capital sets $\mathcal K_z$, a law of motion
$K'=\Psi(z,K,z')$ keeping capital in the sets, labour $L(z,K)$, and value and saving functions of
$(z,K,a,e)$ satisfying the Bellman equation with continuation at $(z',\Psi(z,K,z'))$, together with
consistency of aggregation: each $(z,K)$ carries a nonempty set of distributions with mean $K$ and
labour $L(z,K)$, mapped by the saving policy and the employment transitions into the set at
$(z',\Psi(z,K,z'))$. It \emph{represents} $\mu^*$ if at every history the capital lies in
$\mathcal K_z$, the distribution lies in the set at $(z,K)$, and households save by the $(z,K)$
policy.

\begin{theorem}\label{thm:iff}
For every sequence-state equilibrium, $(z,K)$-sufficiency holds if and only if some recursive
equilibrium on $(z,K)$ represents it. For small capital shares (as in Theorem~\ref{thm:stationary}):
when $(z,K)$-sufficiency holds, the capital sets are compact, the law of motion $\Psi$ is continuous
on them, and the law of $(z,K)$ is the unique invariant distribution of the $(z,K)$ process on the
capital sets, and it is ergodic.
\end{theorem}
\formal{kSufficient\_iff\_exists\_represents, minMPC\_smallShare\_aggregateState, stationary\_aggLaw}

\noindent
Continuity needs no injectivity: capital is a quotient map from the compact set of histories with a
given shock. Uniqueness of the stationary law needs no Lipschitz test functions and no measurable
selection of distributions: bounded continuous test functions are uniformly continuous on the
compact capital sets, and capital forgets old shocks.

\begin{theorem}\label{thm:repagent}
With a single idiosyncratic state, every household holds aggregate capital in the small-share
equilibrium, so capital determines the distribution and $(z,K)$-sufficiency holds: the stochastic
growth model of \citet{brockmirman1972} is recovered as a recursive equilibrium on $(z,K)$ with a
unique stationary distribution.
\end{theorem}
\formal{minMPC\_smallShare\_repAgent}

\paragraph{An order route.} If the cross-sections the equilibrium visits at each shock form a chain
for first-order stochastic dominance in assets, then capital determines the cross-section, because
two comparable cross-sections with the same mean are equal. The chain property follows from two
conditions: the law of motion preserves dominance, and after ordered shocks the next cross-sections
are ordered whatever the pasts. The first holds, by Tarski comparative statics, when saving rises
with the capital path (Theorem~\ref{thm:monotone}); the second reduces to finitely many comparisons
between extreme histories.

\begin{theorem}\label{thm:order}
Dominance with equal means is equality. If saving rises with the capital path, sequential
equilibrium is unique, the equilibrium is continuous in the history, a better newest shock raises
the cross-section, and the extreme comparisons hold, then $(z,K)$-sufficiency holds.
\end{theorem}
\formal{eq\_of\_dominates\_of\_mean\_eq, kInjective\_of\_chain, kSufficient\_of\_savingMonotone,
kSufficient\_of\_extremes}

\noindent
The extreme comparisons say that one period of a better shock outweighs a permanent boom against a
permanent bust. They force capital to encode the whole shock history, which needs the economy to
forget at least half of a capital deviation each period with two shocks. They cannot hold at the
persistence of Krusell and Smith's calibration, which says nothing about $(z,K)$-sufficiency there,
only that this sufficient condition is unavailable.

\section{An explicit economy with heterogeneous households}\label{sec:corner}

\subsection{Corner economies}

Let the employed save up to the cap and the unemployed spend everything. Under the employment
design of \citet{krusellsmith1998} --- the employment rate at shock $z$ is $m(z)$, consistent under
the transitions --- the cross-section after yesterday's shock $y$ and today's $z$ puts mass on two
asset levels, $\abar$ for those employed last period and $0$ for the others, jointly with the current
employment state. Capital is $\abar\, m(y)$. Below the root of a history's tree, capital and labour
depend only on the shocks along the node and today's shock, so yesterday's shock enters prices only
through today's capital: $(z,K)$-sufficiency holds with no injectivity.

The corner policy is optimal when three conditions hold along the price path: labour income covers
the floor plus the corner saving; for the employed, $u'(c)\le\beta\E[R'u'(c')]$ at the cap; for the
unemployed, $u'(c)\ge\beta\E[R'u'(c')]$ at zero, where $c'$ is consumption next period under the
corner policy itself. The proof is a verification argument: the value of following the corner policy
is today's utility plus a term independent of assets; against it the objective is concave, the
tangent lines today and tomorrow bound the gain from leaving the corner by
$(g-a)\,(u'(c)-\beta\E R'u'(c'))$, which has the right sign; so the policy's value solves the Bellman
equation and is the value function. With prices bounded in $[R_{lo},R_{hi}]\times[w_{lo},w_{hi}]$ and
marginal utility decreasing, the conditions hold along every price path once they hold at their
worst cases.
\formal{policy\_eq\_cornerPolicy, policyValue\_eq\_valueFunction, cornerConditions\_of\_bounds}

\subsection{The economy}

\begin{center}
\begin{tabular}{ll}
\toprule
Aggregate shocks & two, i.i.d.\ with probability $1/2$ \\
Employment & efficiency $2$ (employed), $1/4$ (unemployed) \\
Employment rates & $4/5$ (bad), $17/20$ (good); job finding $49/50$ \\
Job loss & $49/200,\ 73/400,\ 197/850,\ 147/850$ by (yesterday, today) \\
Preferences & logarithmic, $\beta=24/25$; asset cap $\abar=1/16$ \\
Technology & Cobb--Douglas, $\alpha=1/10$, TFP $1/4$, $\delta=1/10$ \\
Boxes & $K^{1/10}\in[11/15,\,3/4]$, $L^{1/10}\in[21/20,\,16/15]$ \\
\bottomrule
\end{tabular}
\end{center}

\noindent
The box endpoints are tenth powers of rationals, so the powers $1/10$ and $9/10$ in the prices are
exact and every inequality is checked in exact rational arithmetic. Gross returns lie in about
$[1.42,1.63]$ and wages in about $[0.155,0.161]$. Capital is $1/20$ after a bad shock and $17/320$
after a good one, labour $33/20$ or $139/80$: both inside the boxes at every node of the equilibrium
price tree, so prices are the true Cobb--Douglas marginal products.

\begin{theorem}\label{thm:corner}
In this economy: sequential equilibrium is unique from every initial state; the sequence-state
equilibrium exists, is unique, and is the corner one; $(z,K)$-sufficiency holds and a recursive
equilibrium on $(z,K)$ represents it; and its law of motion on $(z,K)$ has a unique stationary
distribution on the capital sets --- weight $1/4$ on each pair of today's shock and capital from
yesterday's shock --- which is ergodic.
\end{theorem}
\formal{cornerUnique, cornerWitness, cornerStationary, ergodic\_statLaw, seqPrices\_mem\_box}

\noindent
Uniqueness comes cheaply because the corner conditions hold along every price path: the household's
policy is the corner one whatever it expects, so the distribution path does not depend on prices.
Three features limit the example. Assets sit at corners rather than at interior precautionary
levels. The gross return exceeds $1/\beta$, which is what makes the cap bind for the employed. And
with capital share $1/2$ and labour inside its box the corner economy cannot exist: the employed's
net income over the unemployed's spending is bounded by roughly $(1-\alpha)/\alpha$, too small for
both corner conditions.

\section{An economy with interior saving}\label{sec:interior}

The corner economy keeps the cross-section on two asset levels, so its history is short. Here the
cross-section carries its whole history, saving is interior, and $(z,K)$ is still sufficient,
because capital encodes the history. The route is the order route of Section~\ref{sec:aggregate},
with an economy in which its hypotheses can be checked.

\subsection{Separation from ordered cash}

Three conditions make the two hard hypotheses of Theorem~\ref{thm:order} free.
\begin{itemize}[leftmargin=*,itemsep=2pt]
\item \textbf{Prices depend on the shock only}, as with a linear technology $Y=A(z)K+B(z)L$, whose
  marginal products are exact. Then the household's problem does not depend on the path of
  capital, so saving rises (weakly) with capital, and sequential equilibrium is unique: the
  distribution path does not depend on prices.
\item \textbf{Shocks are i.i.d.}: next period's aggregate shock and employment are drawn
  independently of today's. Then the value of next period's assets is the same function at the
  root of every history, so today's saving is one nondecreasing function of cash on hand (a
  Topkis argument across histories).
\item \textbf{Cash is ordered by the shock}: after a worse shock every household has less cash than
  every household after a better one.
\end{itemize}
Under these, every household saves less after the worse shock, so the next cross-sections are
ordered for any two current ones. This is the separation condition of the order route, directly,
without the extreme histories.

\begin{theorem}\label{thm:orderedcash}
If prices depend on the shock only, shocks are i.i.d.\ and cash on hand is ordered by the shock,
then in every sequence-state equilibrium $(z,K)$ is a sufficient state.
\end{theorem}
\formal{kSufficient\_of\_orderedCash, sepOnHistories\_of\_orderedCash, policy\_le\_policy\_of\_cont,
cont\_root\_eq, uniqueSequentialEquilibrium\_of\_exogenous}

\subsection{Interior saving}

Interiority is proved by two corner deviations. Extra assets $\delta$ are worth at least
$R\delta/(\text{cash}+R\delta)$ next period: keep the saving of the poorer state and consume the
extra return, and use the tangent line of the logarithm. The last $\delta$ below $b$ is worth at
most $R\delta/(\text{cash}(b)-\abar-R\delta)$: the poorer state copies the richer one's saving. So
saving is positive at cash $x$ if
\[
  \frac{\delta}{x-\delta} < \beta\,\E\!\left[\frac{R'\delta}{w'e'+R'\delta}\right],
\]
and below the cap if
\[
  \frac{\delta}{x-\abar+\delta} > \beta\,\E\!\left[\frac{R'\delta}{(R'-1)\abar+w'e'-R'\delta}\right].
\]
\formal{value\_sub\_ge, value\_sub\_le, policy\_pos\_of, policy\_lt\_cap\_of}

\subsection{The economy}

\begin{center}
\begin{tabular}{ll}
\toprule
Aggregate shocks & two, i.i.d.\ with probability $1/2$ \\
Efficiency & $1$ or $3/4$, i.i.d.\ with probability $1/2$ \\
Technology & linear: gross return $6/5$; wage $1$ (good), $1/3$ (bad) \\
Preferences & logarithmic, $\beta=1/2$; asset cap $\abar=25/72$; floor $1/100$ \\
\bottomrule
\end{tabular}
\end{center}

\noindent
After a bad shock cash is at most $\tfrac65\cdot\tfrac{25}{72}+\tfrac13=\tfrac34$; after a good
one it is at least $\tfrac34$. With $\delta=1/3000$, for low-efficiency households after a good
shock (cash in $[3/4,\,7/6]$) the first condition reads $1/2249 < \tfrac18(1/2501+1/1876+3/2503+1/626)$
and the second $3/7378 > \tfrac18\cdot18\,(1/48107+1/36857+1/18107+1/14357)$, both in exact
rational arithmetic.

\begin{theorem}\label{thm:interior}
This economy has a sequence-state equilibrium. In every one, $(z,K)$ is a sufficient state, and
after every good shock at least half of the households --- all those with low efficiency, whatever
their assets --- save strictly between $0$ and $\abar$.
\end{theorem}
\formal{InteriorWitness.exists\_interiorWitness, InteriorWitness.interiorWitness,
InteriorWitness.saving\_interior, InteriorWitness.orderedCash}

\noindent
The costs are plain. Prices do not respond to capital, so the equilibrium feedback that makes
Krusell--Smith hard is absent; what remains is the aggregation question itself, whether aggregate
capital pins down a cross-section that keeps its whole history, and here it does. Shocks are
i.i.d.\ rather than persistent. And the household is impatient: with these crude interiority
bounds the windows are tiny for $\beta\ge 0.8$, so the period should be read as long. We do not
prove that interior saving is strictly increasing in assets, which would make the cross-section
infinitely supported.

\section{What is open}\label{sec:open}

\begin{enumerate}[leftmargin=*,itemsep=2pt]
\item \textbf{$(z,K)$ at interior, persistent calibrations.} The heterogeneous example of
  Section~\ref{sec:corner} has assets at corners. Individual nonlinearity of saving does not by
  itself defeat $(z,K)$-sufficiency --- only the cross-sections the equilibrium visits matter --- so
  the question is open in both directions. The order route of Section~\ref{sec:aggregate} needs fast
  forgetting; an economy with interior saving, i.i.d.\ wage shocks and complementarity is the next
  candidate.
\item \textbf{Approximate aggregation.} A quantitative version of $(z,K)$-sufficiency, with errors
  bounded by how far apart visited cross-sections with the same $(z,K)$ can be. In the small-share
  class the bound is of order $\alpha$, but so is the error of forecasting with less information.
\item \textbf{Measurable equilibria beyond small shares.} Stationarity directly from the sequential
  equilibrium needs the sequence-state equilibrium to be measurable in the history; outside the
  small-share class the fixed point gives no such guarantee.
\item \textbf{Uniqueness at realistic capital shares.} Theorem~\ref{thm:smallshare} is qualitative;
  sharper bounds on the consumption floor and on curvature would be needed for a quantitative
  threshold.
\item \textbf{Full spotlessness}, open also in \citet{cao2020}; and the asset cap, a constraint
  throughout.
\end{enumerate}


\bibliographystyle{plainnat}
\bibliography{refs}

\appendix
\section{Index of formal statements}\label{app:index}

All names are in the namespace \lean{LeanEconomics.KrusellSmith} (most in \lean{KSModel}); files
are in \texttt{LeanEconomics/KrusellSmith/} unless noted.

{\footnotesize
\begin{longtable}{>{\raggedright}p{0.27\textwidth}>{\raggedright}p{0.38\textwidth}>{\raggedright\arraybackslash}p{0.25\textwidth}}
\toprule
Result & Formal name & File \\
\midrule
\endhead
Existence (Thm~\ref{thm:existence}) & \lean{exists\_sequentialEquilibrium} & \texttt{Equilibrium} \\
Labour consistency & \lean{exists\_sequentialEquilibrium\_labour} & \texttt{Labour} \\
Log, CRRA (Thm~\ref{thm:existenceInada}) & \lean{KSInada.exists\_sequentialEquilibrium} & \texttt{Inada} \\
Cobb--Douglas (Thm~\ref{thm:existenceCD}) & \lean{logEconomy\_exists\_equilibrium} & \texttt{CobbDouglas} \\
Clamp never binds & \lean{crra\_exists\_unclamped\_equilibrium} & \texttt{CobbDouglasCapital} \\
Minimal state under uniqueness (Thm~\ref{thm:minimal}) & \lean{KSModel.isRecursive} & \texttt{Recursive} \\
Value-function state (Thm~\ref{thm:valuestate}) & \lean{recursiveEquilibrium\_valueState} & \texttt{ValueStateExistence} \\
History state (Thm~\ref{thm:history}) & \lean{exists\_seqStateEquilibrium} & \texttt{SequenceState\allowbreak Equilibrium} \\
DGMM stationary (Thm~\ref{thm:dgmm}) & \lean{exists\_stationaryRecursiveEquilibrium} & \texttt{StationaryEquilibrium} \\
Conditionally spotless & \lean{exists\_conditionallySpotlessEquilibrium} & \texttt{SpotlessEquilibrium} \\
Monotone route (Thm~\ref{thm:monotone}) & \lean{SavingMonotone.isRecursive} & \texttt{GreatestEquilibrium} \\
Uniqueness, CRRA (Thm~\ref{thm:smallshare}) & \lean{crra\_existsUnique\_sequentialEquilibrium} & \texttt{CRRAContraction} \\
Small capital shares & \lean{minMPC\_smallShare} & \texttt{MinMPCSmallShare} \\
Unique sequence state (Thm~\ref{thm:sequnique}) & \lean{minor\_minMPC\_smallShare\_existsUnique\_seqStateEquilibrium} & \texttt{DoeblinUniqueness} \\
Matching (Thm~\ref{thm:matching}) & \lean{minMPC\_smallShare\_matching} & \texttt{HistoryMarginalsCD} \\
Factor theorem (Thm~\ref{thm:factor}) & \lean{stateLaw\_tomorrow}, \lean{isInvariant\_stateLaw} & \texttt{StationarySequence} \\
Unique stationary (Thm~\ref{thm:stationary}) & \lean{minMPC\_smallShare\_unique\_stationary} & \texttt{StationaryUnique} \\
$(z,K)$ iff (Thm~\ref{thm:iff}) & \lean{kSufficient\_iff\_exists\_represents} & \texttt{AggregateStateIff} \\
$(z,K)$, small shares & \lean{minMPC\_smallShare\_aggregateState} & \texttt{AggregateStateIff} \\
Representative agent (Thm~\ref{thm:repagent}) & \lean{minMPC\_smallShare\_repAgent} & \texttt{AggregateStateRepAgent} \\
Order route (Thm~\ref{thm:order}) & \lean{kSufficient\_of\_extremes} & \texttt{OrderExtremes} \\
Corner verification & \lean{policy\_eq\_cornerPolicy} & \texttt{CornerVerification} \\
Policy evaluation & \lean{policyValue\_eq\_valueFunction} & \texttt{DynamicProgramming/\allowbreak PolicyEvaluation} \\
Corner conditions from bounds & \lean{cornerConditions\_of\_bounds} & \texttt{CornerPrimitives} \\
Corner economy (Thm~\ref{thm:corner}) & \lean{CornerWitness.cornerUnique} & \texttt{CornerUnique} \\
True marginal products & \lean{CornerWitness.seqPrices\_mem\_box} & \texttt{CornerWitness} \\
Stationary law, ergodic & \lean{CornerWitness.cornerStationary}, \lean{CornerWitness.ergodic\_statLaw} & \texttt{CornerStationary} \\
\bottomrule
\end{longtable}
}


\end{document}
